\documentclass[11pt,a4paper]{article}
\usepackage[utf8]{inputenc}
\usepackage[T1]{fontenc}
\usepackage{amsmath,amssymb}
\usepackage{booktabs}
\usepackage{hyperref}
\usepackage{geometry}
\geometry{margin=2.5cm}

\title{Assignment 1 -- Exercises 1 \& 2\\\large FAECTOR Excellence Program\\Energy System Optimization and Pricing}
\author{}
\date{}

\begin{document}
\maketitle

\section*{Do the solutions answer Q1 and Q2?}

Yes. The implementation in \texttt{Code/Script\_Ex1\&2.py} addresses both questions:
\begin{itemize}
  \item \textbf{Question 1:} The market-clearing optimization is formulated and solved using the provided data (\texttt{Data/Realistic Data Generators (1).xlsx} and \texttt{Data/Realistic Data Demand.xlsx}).
  \item \textbf{Question 2:} The price in each time period is taken as the dual (shadow price) of the power-balance constraint; the script also supports varying $\alpha$ and $\lambda$ to observe how these prices change.
\end{itemize}

\section{Question 1: Market-clearing formulation and code}

\subsection{Data}
\begin{itemize}
  \item \textbf{Generators:} 8 units with maximum capacity $P_g$ and variable cost $MC_g$ (constant over time).
  \item \textbf{Demand:} 168 time periods and 5 renewable scenarios; one scenario is selected for each run.
\end{itemize}

\subsection{Decision variables}
\begin{itemize}
  \item $p_{g,t} \geq 0$: conventional generation from generator $g$ in period $t$.
  \item $c_t \geq 0$: curtailed renewable energy in period $t$.
\end{itemize}

\subsection{Objective}
Minimize total cost (conventional generation cost plus curtailment penalty):
\begin{equation}
  \min \sum_{g,t} MC_g \cdot p_{g,t} + \lambda \sum_t c_t.
\end{equation}

\subsection{Constraints}
\begin{itemize}
  \item \textbf{Power balance} (equality, no free disposal): for each $t$,
    \begin{equation}
      \sum_g p_{g,t} + \bigl(\alpha \cdot R_t - c_t\bigr) = D_t,
    \end{equation}
    where $R_t$ is the renewable production in the chosen scenario and $D_t$ is demand.
  \item \textbf{Generation limits:} $0 \leq p_{g,t} \leq P_g$ for all $g,t$.
  \item \textbf{Curtailment limits:} $0 \leq c_t \leq \alpha R_t$ for all $t$.
\end{itemize}

Demand is perfectly inelastic; $\alpha \in [0,1]$ scales renewable availability and $\lambda > 0$ is the curtailment penalty. The problem is implemented as a linear program and solved with \texttt{scipy.optimize.linprog} (HiGHS).

\section{Question 2: Prices and sensitivity to $\alpha$ and $\lambda$}

\subsection{Price definition}
The \emph{market price} in period $t$ is the dual variable (shadow price) of the power-balance constraint for that period. In the solution object these are \texttt{result.eqlin.marginals}; they are reported in the balance table and in \texttt{exercise2\_prices.csv}.

\subsection{Observations}
\begin{itemize}
  \item \textbf{Effect of $\alpha$:} For low $\alpha$ (e.g.\ 0.1 or 0.5), renewable supply is small relative to demand, so there is no curtailment and prices remain positive (between the lowest and highest $MC_g$). For $\alpha = 1$, renewable availability is high; in some periods surplus is curtailed and the dual can be \emph{negative}, reflecting that an extra unit of demand would allow using less costly curtailed energy (reducing the penalty $\lambda c_t$).
  \item \textbf{Effect of $\lambda$:} When there is curtailment, a higher $\lambda$ increases the cost of curtailing; the shadow price of the balance constraint (the market price) becomes more negative in those periods. So $\lambda$ directly influences the magnitude of negative prices when curtailment binds.
\end{itemize}

The script options \texttt{--compare-alphas} and \texttt{--compare-lambdas} run the model on a grid of values and write summary statistics (average/min/max price, number of periods with negative price, total curtailment) to \texttt{exercise2\_price\_sensitivity.csv}.

\section*{Code and outputs}
\begin{itemize}
  \item Script: \texttt{Code/Script\_Ex1\&2.py}.
  \item Default data: \texttt{Data/Realistic Data Demand.xlsx}, \texttt{Data/Realistic Data Generators (1).xlsx}.
  \item With \texttt{--save-results}, outputs are written to \texttt{Data/}: \texttt{exercise1\_dispatch.csv}, \texttt{exercise1\_balance.csv}, \texttt{exercise2\_prices.csv}, and optionally \texttt{exercise2\_price\_sensitivity.csv}.
\end{itemize}

\end{document}
